%% maestro2tex -- generated table; do not edit by hand.
\begin{table}[htbp]
  \centering
  \caption[{Finite-loop-gain error against feedback resistance at the typical statistical process point,\,\dots}]{Finite-loop-gain error against feedback resistance at the typical statistical process point, \SI{37}{\celsius}, truth current swept \SIrange{-2.3}{2.3}{\micro\ampere} with $R_{\mathrm{ct}} = \SI{10}{\kilo\ohm}$ held fixed across the whole sweep. \textbf{The 315 and \SI{560}{\kilo\ohm} rows are NOT the accuracy of those gain taps} -- $R_{\mathrm{ct}} = \SI{10}{\kilo\ohm}$ is a stiff, non-operational electrode at those gains (the operational electrode is \SI{1}{\mega\ohm}, Table~\ref{tab:px-hg1m-acc}). They demonstrate that the residual error is finite closed-loop gain alone: the zero-crossing slope error is $-1/T$ to within \SI{4}{\percent} at every tap, and the worst-case error scales with $R_f$ as $T$ falls. The LSB column is referred to the nominal \SI{2.441}{\milli\volt} step; the converter's measured step is \SI{1.591}{\milli\volt} (Table~\ref{tab:adc-transfer}), so multiply it by \num{1.535} to read it as converter codes.}
  \label{tab:px-hg-demo}
  \begin{tabular}{lrrr}
    \toprule
    $R_f$ & Worst error over the sweep~[LSB] & Zero-crossing error~[\si{\nano\ampere}] & Zero-crossing slope error ($=-1/T$) \\
    \midrule
    35000 & 0.271 & $-17.704$ & $-0.00052$ \\
    140000 & 0.407 & $-4.426$ & $-0.00118$ \\
    315000 & 0.935 & $-1.967$ & $-0.00233$ \\
    560000 & 3.101 & $-1.107$ & $-0.00394$ \\
    \bottomrule
  \end{tabular}
\end{table}
